\chapter{Introduction}

\section{États de l'art des modèles océaniques}
Depuis le début du XX\textsuperscript{e} siècle, les scientifiques s'attellent à la prédiction météorologique à long terme, d'abord par des modèles théoriques dès 1922 avec Lewis Fry Richardson qui propose l'idée de la prévision météorologique numérique \cite{richardson_weather_1922}. Cette idée se concrétise dès le début des années 1950 avec l'ordinateur ENIAC et l'équipe dirigée par John von Neumann et Jule Charney \cite{charney_numerical_1950}. 

Concernant les modèles d'océanographie, il faut attendre le milieu du XX\textsuperscript{e} siècle avec Henri Stommel et son modèle dynamique du Gulf Stream \cite{stommel_westward_1948}. Ce modèle a ensuite été affiné quelques années plus tard par Munk \cite{munk_wind-driven_1950} en prenant en compte la dissipation ainsi que l'interaction avec les vents, ce qui a permis de mieux comprendre les gyres océaniques à grande échelle. Le développement de modèles numériques tridimensionnels d'océanographie commence dans les années 1960 avec Kirk Bryan et Michael Cox \cite{bryan_nonlinear_1968, bryan_nonlinear_1968-1}, qui ont créé \textbf{MOM}, un modèle encore utilisé actuellement pour étudier le système climatique océanique mondial. À ce jour, et depuis les années 1990, on utilise des modèles couplés océan-atmosphère. L'un des premiers est \textbf{OASIS}, un coupleur océan-atmosphère autonome \cite{terray_oasis_1995} qui a permis une simulation de 10 ans de l'océan Pacifique. Depuis, plusieurs modèles couplés ont émergé, comme \textbf{LIM3}, un modèle océan-banquise \cite{vancoppenolle_lim3_2008}, le modèle de circulation côtière \textbf{ROMS} \cite{moore_regional_2011} et également \textbf{NEMO} \cite{madec_ocean_1997}. Récemment, le développement de modèles océaniques couplés avec un système de \textit{machine learning} adapté aux sciences de la Terre s'accélère \cite{wang_coupled_2024}, ce qui permettra, dans le futur, de réaliser des simulations plus rapidement avec moins d'erreurs.

\section{Émergence des modèles probabilistes}
Les modèles évoqués précédemment sont des modèles déterministes, qui ne décrivent pas de manière cohérente les incertitudes du modèles pour les systèmes d'assimilation de données. Ces incertitudes ne sont pas les incertitudes d'ordre matériel comme la précision des machines ou les schémas numériques, mais bien les incertitudes liées à la dynamique, comme la turbulence non résolue, le forçage ou encore les conditions initiales.. Afin de pouvoir considérer ces modèles comme des modèles prédictifs, selon les définitions données dans \cite{brier_verification_1950, toth_7_2003}, il faut tenir compte de ces incertitudes dans le modèle en lui-même.

Dès le début du XXI\textsuperscript{e} siècle, la modélisation de ces incertitudes devient un enjeu central. Par exemple, \cite{buizza_stochastic_1999} analyse l’impact des paramétrisations randomisées sur l’ECMWF EPS, un modèle météorologique.. Concernant l'océanographie, Jean-Michel Brankart a développé des méthodes de paramétrisations stochastiques appliquées au modèle \textbf{NEMO} \cite{brankart_generic_2015}, méthodes qui seront adaptées pour produire des ensembles de simulations avec \textbf{CROCO}, un modèle se distinguant de \textbf{NEMO} par sa coordonnée verticale, qui varie selon la bathymétrie et permet une meilleure résolution des phénomènes côtiers.